\section{The new gauge group and its covariant derivative}
\label{sec:gauge}

\subsection{The symmetry that the pair selects}

Let a $2\times2$ complex matrix $\Upsilon_{U2}$ act on the doublet index,
$\Psi_{U2}\to(\Upsilon_{U2}\otimes\mathsf 1_{16})\Psi_{U2}$; such a transformation commutes with
Spin(7,1), because it does not touch the spinor index. (Below, a matrix that acts on one of the two
indices only is written without the unit matrix of the other: $\gamma^A_{1671}$ acting on
$\Psi_{U2}$ means $\mathsf 1_2\otimes\gamma^A_{1671}$, and $\Upsilon_{U2}$ acting on $\Psi_{U2}$ means
$\Upsilon_{U2}\otimes\mathsf 1_{16}$.) The kinetic and mass terms of the doublet,
$\bar\Psi_{U2}(i\gamma^A_{1671}\partial_A-m_{U2})\Psi_{U2}$, are the sum of the same terms for
$\psi^1_{U2}$ and $\psi^2_{U2}$ with the common mass $m_{U2}$, and $\bar\Psi_{U2}=
\Psi^\dagger_{U2}\beta_{1671}$. Under a constant $\Upsilon_{U2}$ they become
$\Psi^\dagger_{U2}\Upsilon^\dagger_{U2}\beta_{1671}(i\gamma^A_{1671}\partial_A-m_{U2})\Upsilon_{U2}\Psi_{U2}$,
which equals the original for all fields if and only if
$\Upsilon^\dagger_{U2}\Upsilon_{U2}=\mathsf 1_2$: the transformation must be unitary. For every $2\times2$ matrix the entries of $\Upsilon^T_{U2}\mathsf E\,\Upsilon_{U2}$ are its
$2\times2$ minors, so that
\begin{equation}
\Upsilon^T_{U2}\,\mathsf E\,\Upsilon_{U2}=\det(\Upsilon_{U2})\,\mathsf E,
\qquad
\Phi^\pm_{U2}\to\det(\Upsilon_{U2})\,\Phi^\pm_{U2},
\label{eq:detE}
\end{equation}
by Eq.~\eqref{eq:pairs} (R3-A9). A unitary transformation leaves the spin-zero pair invariant if
and only if its determinant is 1: the pair selects $SU(2)$. Together with the overall phase,
which counts the fermions, the unitary transformations form
\begin{equation}
\begin{aligned}
U(2)_{U2}&=\big(SU(2)_{U2}\times U(1)_{U2}\big)/\mathbb Z_2,\\
\Upsilon_{U2}&=\exp\!\big(\tfrac i2\boldsymbol\Theta_{U2}\cdot\boldsymbol\tau_{U2}\big)\exp\!\big(iy_{U2}\Theta^0_{U2}\big),
\end{aligned}
\label{eq:group}
\end{equation}
with the Pauli matrices $\boldsymbol\tau_{U2}=(\tau^1_{U2},\tau^2_{U2},\tau^3_{U2})$ acting on the
doublet index, three real parameters $\boldsymbol\Theta_{U2}$, one real parameter
$\Theta^0_{U2}$, and the $U(1)_{U2}$ charge $y_{U2}=\tfrac12$ of the doublet. The quotient by
$\mathbb Z_2$ removes the double counting of the element $-\mathsf 1_2$, which belongs to
$SU(2)_{U2}$ and equals $\exp(iy_{U2}\Theta^0_{U2})$ at $\Theta^0_{U2}=2\pi$. The determinant is
$\det\Upsilon_{U2}=\exp(2iy_{U2}\Theta^0_{U2})=\exp(i\Theta^0_{U2})$. The value $y_{U2}=\tfrac12$
copies the hypercharge of the Higgs doublet, Eq.~\eqref{eq:sm}; with it the pair has charge 1.
We make this symmetry local, $\boldsymbol\Theta_{U2}(X)$ and $\Theta^0_{U2}(X)$, and call it the
new gauge group. The group is the group of the electroweak theory; what is new is where it acts
(on an internal index of 16-component fermions of the $7+1$ spacetime), why it is there (it is
the symmetry of the only spin-zero pair, Theorem~1), and which field breaks it (a pair of
fermions, not an elementary scalar).

\subsection{The covariant derivative}

Modelled on Eq.~\eqref{eq:sm}, the new covariant derivative is
\begin{equation}
\begin{aligned}
D_{U2A}\Psi_{U2}&=\Big(\partial_A+\tfrac i2\,g_{U2}\,\boldsymbol\tau_{U2}\cdot\boldsymbol W_{U2A}\\
&\qquad+i\,g'_{U2}\,y_{U2}\,B_{U2A}\Big)\Psi_{U2}\\
&\equiv(\partial_A+i\mathcal W_{U2A})\Psi_{U2},
\end{aligned}
\label{eq:D}
\end{equation}
with a triple of real vector fields $\boldsymbol W_{U2A}=(W^1_{U2A},W^2_{U2A},W^3_{U2A})$ for
$SU(2)_{U2}$, a real vector field $B_{U2A}$ for $U(1)_{U2}$, the couplings $g_{U2}$ and $g'_{U2}$,
and the Hermitian $2\times2$ matrix $\mathcal W_{U2A}=\tfrac12g_{U2}\boldsymbol\tau_{U2}\cdot\boldsymbol
W_{U2A}+g'_{U2}y_{U2}B_{U2A}\mathsf 1_2$. It acts on the doublet index, on each of the 16 spinor
components alike. Under a local transformation $\Psi'_{U2}=\Upsilon_{U2}\Psi_{U2}$ the connection
transforms as
\begin{equation}
\mathcal W'_{U2A}=\Upsilon_{U2}\mathcal W_{U2A}\Upsilon^{-1}_{U2}+i(\partial_A\Upsilon_{U2})\Upsilon^{-1}_{U2},
\label{eq:Wtrans}
\end{equation}
and then $D'_{U2A}\Psi'_{U2}=\Upsilon_{U2}D_{U2A}\Psi_{U2}$: indeed
$(\partial_A+i\mathcal W'_{U2A})\Upsilon_{U2}\Psi_{U2}=(\partial_A\Upsilon_{U2})\Psi_{U2}+\Upsilon_{U2}\partial_A\Psi_{U2}
+i\Upsilon_{U2}\mathcal W_{U2A}\Psi_{U2}-(\partial_A\Upsilon_{U2})\Psi_{U2}$ (R3-G1). The new connection
$\mathcal W'_{U2A}$ is again Hermitian, because $\Upsilon_{U2}$ is unitary, and its $U(1)_{U2}$ part
is $B'_{U2A}=B_{U2A}-\partial_A\Theta^0_{U2}/g'_{U2}$ (R3-G2). To first order in
$\boldsymbol\Theta_{U2}$ the $SU(2)_{U2}$ part reads
$\boldsymbol W'_{U2A}=\boldsymbol W_{U2A}-\boldsymbol\Theta_{U2}\times\boldsymbol W_{U2A}
-\partial_A\boldsymbol\Theta_{U2}/g_{U2}$ ($\times$ is the cross product of triples); it follows from
$\tau^1_{U2}\tau^2_{U2}-\tau^2_{U2}\tau^1_{U2}=2i\tau^3_{U2}$ and the two cyclic permutations (R3-A9),
which turn the commutator of two combinations of Pauli matrices into the combination with the
cross product of the coefficients.

The commutator of two covariant derivatives defines the field strength,
\begin{equation}
\begin{aligned}
[D_{U2A},D_{U2C}]\Psi_{U2}&=i\mathcal F_{U2AC}\Psi_{U2},\\
\mathcal F_{U2AC}&=\partial_A\mathcal W_{U2C}-\partial_C\mathcal W_{U2A}\\
&\quad+i[\mathcal W_{U2A},\mathcal W_{U2C}]\\
&=\tfrac12g_{U2}\boldsymbol\tau_{U2}\cdot\boldsymbol W_{U2AC}\\
&\quad+g'_{U2}y_{U2}B_{U2AC}\mathsf 1_2,
\end{aligned}
\label{eq:F}
\end{equation}
with $\boldsymbol W_{U2AC}=\partial_A\boldsymbol W_{U2C}-\partial_C\boldsymbol W_{U2A}-g_{U2}\boldsymbol
W_{U2A}\times\boldsymbol W_{U2C}$ and $B_{U2AC}=\partial_AB_{U2C}-\partial_CB_{U2A}$ (R3-G3, R3-G4).
It transforms by conjugation, $\mathcal F'_{U2AC}=\Upsilon_{U2}\mathcal F_{U2AC}\Upsilon^{-1}_{U2}$
(R3-G5), so the traces $\mathrm{Tr}(\mathcal F_{U2AC}\mathcal F^{AC}_{U2})=\tfrac12g_{U2}^2
\boldsymbol W_{U2AC}\cdot\boldsymbol W^{AC}_{U2}+2g'^2_{U2}y_{U2}^2B_{U2AC}B^{AC}_{U2}$ and
$\mathrm{Tr}\,\mathcal F_{U2AC}=2g'_{U2}y_{U2}B_{U2AC}$ are invariant, and so is the Lagrangian of
the gauge fields, $-\tfrac14\boldsymbol W_{U2AC}\cdot\boldsymbol W^{AC}_{U2}-\tfrac14B_{U2AC}B^{AC}_{U2}$.

On the eight-dimensional spacetime the action is dimensionless when the Dirac field has the
dimension mass$^{7/2}$ and the gauge fields mass$^{3}$; the couplings $g_{U2}$ and $g'_{U2}$ then
have the dimension mass$^{-2}$.

\subsection{The charges of the pairs}

By Eq.~\eqref{eq:detE} the pairs are $SU(2)_{U2}$ singlets and carry the $U(1)_{U2}$ charge
$2y_{U2}=1$: $\Phi^\pm_{U2}\to\exp(i\Theta^0_{U2})\Phi^\pm_{U2}$ (R3-G6). Their covariant derivative
contains only $B_{U2A}$. Since $(\tau^1_{U2})^T\mathsf E+\mathsf E\tau^1_{U2}=0$, and the same for
$\tau^2_{U2}$ and $\tau^3_{U2}$ (R3-A9), $\mathcal W^T_{U2A}\mathsf E+\mathsf E\mathcal W_{U2A}
=2g'_{U2}y_{U2}B_{U2A}\mathsf E$, and the product rule gives
\begin{equation}
\begin{aligned}
(\partial_A+ig'_{U2}B_{U2A})\Phi^\pm_{U2}={}&\tfrac12(D_{U2A}\Psi_{U2})^T(\mathsf E\otimes C^\pm_{1671})\Psi_{U2}\\
&+\tfrac12\Psi^T_{U2}(\mathsf E\otimes C^\pm_{1671})D_{U2A}\Psi_{U2}.
\end{aligned}
\label{eq:Dpair}
\end{equation}
A condensate of $\Phi^+_{U2}$ is therefore invariant under $SU(2)_{U2}$ and changes its phase under
$U(1)_{U2}$: if it forms, it breaks $U(2)_{U2}\to SU(2)_{U2}$, and only $B_{U2A}$ can acquire a mass
(Sec.~\ref{sec:breaking}). In the electroweak theory the Higgs doublet carries both kinds of
charge, and the boson that becomes massive is a mixture of the third component of $\boldsymbol W_{\rm SM}$
and of $B_{\rm SM}$; here the order parameter
is an $SU(2)_{U2}$ singlet, and there is no mixing.

\subsection{Chirality and anomalies}

In eight dimensions the gauge anomaly of Weyl fermions comes from one-loop diagrams with five
external gauge bosons; it is governed by the fifth power of the field strength, and the mixed
gauge-gravitational anomalies by its third and first powers multiplied by curvature
terms~\cite{AlvarezGaumeWitten1984}. For a $U(1)$ the three coefficients are the sums of the fifth, of the third and of the first
powers of the charges of the left-handed Weyl components.
For $SU(2)_{U2}$ alone the symmetrized traces of three and of five Pauli matrices vanish
(R3-B2): the group has no perturbative anomaly of its own.

A Dirac doublet consists of a left-handed and a right-handed doublet. Because complex conjugation
exchanges the chiralities on the $7+1$ spacetime (Sec.~\ref{sec:spacetime}), the right-handed
doublet of charge $y_{U2}$ is equivalent to a left-handed doublet in the conjugate representation,
of charge $-y_{U2}$, and the three coefficients vanish (R3-B2). More strongly, the Dirac doublet
admits the gauge-invariant mass term $m_{U2}\bar\Psi_{U2}\Psi_{U2}$; a regulator field of large
mass of the same kind (Pauli and Villars) then respects $U(2)_{U2}$, so every gauge anomaly of the
Dirac doublet cancels. A doublet of one chirality, by contrast, has the coefficients
$(\tfrac1{16},\tfrac14,1)$ (two components of charge $\tfrac12$; R3-B2): its $U(1)_{U2}$ cannot be
gauged. Since the spin-zero pair joins equal chiralities (Sec.~\ref{sec:pairs}), a chiral doublet
could still pair, but it could neither carry the gauged $U(1)_{U2}$ nor have a mass. The gauged
doublet of this paper is the Dirac doublet; the chiral doublet is followed in
Sec.~\ref{sec:chiral} only to compare its potential.
